\documentclass[11pt]{article}
\usepackage[margin=1in]{geometry}
\usepackage{times}
\usepackage{hyperref}
\hypersetup{colorlinks=true, linkcolor=black, urlcolor=blue, citecolor=black}
\title{A Memo Is Not a Measurement: Wilbert Smith and Project Magnet}
\author{Flyxion \\ \small Independent Researcher}
\date{}
\begin{document}
\maketitle

\begin{abstract}
Wilbert Smith, a Canadian Department of Transport radio engineer, initiated Project Magnet in 1950, a government-sanctioned investigation into geomagnetic effects that Smith privately hoped would lead to an antigravity propulsion breakthrough, and wrote an internal 1950 memorandum, later declassified, describing purported high-level American conversations characterizing flying saucer technology as "the most highly classified subject in the United States government, rating higher even than the hydrogen bomb." This paper argues that Smith's memo documents his own secondhand account of a conversation, not an independently verified government finding, that Project Magnet's actual multi-year experimental program, including a dedicated observatory built specifically to detect anomalous magnetic disturbances associated with claimed unidentified aerial phenomena, concluded without confirming Smith's hoped-for antigravity effect or detecting the anomalies it was built to find, and that the memo's continued circulation as evidence of confirmed American antigravity research rests on treating a single individual's account of a private conversation as though it were independently corroborated government testimony, which it has never been shown to be.
\end{abstract}

\section{Introduction}

Wilbert B. Smith, a senior radio engineer with the Canadian Department of Transport, became interested in the physics of unidentified flying object reports in 1950 and, after what he described as an informal conversation during a Washington visit, wrote an internal memorandum to his superiors stating that he had learned flying saucers were real, were the subject of the most highly classified research program in the United States government, and that research on the underlying propulsion principle was being conducted, a memo that led to Canadian government approval of Project Magnet, a multi-year program under Smith's direction using a dedicated geomagnetic observatory at Shirley Bay, Ontario, intended to detect and study magnetic anomalies Smith believed might be associated both with unidentified aerial phenomena and with a genuine underlying antigravity-adjacent propulsion physics he hoped his own theoretical work could help formalize.

\section{What the Memo Actually Documents}

The 1950 memo, since declassified and widely circulated, records Smith's own account of an informal conversation, not a transcript, recording, or independently corroborated briefing, and names no specific American official as its source in a way that permitted subsequent independent verification of the claimed characterization; no American government document from the same period or subsequently declassified has been identified that corroborates the specific characterization Smith attributed to his unnamed American contacts, leaving the memo's content resting entirely on Smith's own secondhand and unverified account of what he was told, a considerably weaker evidentiary status than its frequent citation as documentary proof of confirmed high-level American antigravity research would suggest.

\section{What Project Magnet Actually Found}

Project Magnet's Shirley Bay observatory operated for several years using magnetometers and other instrumentation specifically intended to detect magnetic disturbances that might correlate with unidentified aerial phenomena sightings or reveal a novel geomagnetic-based propulsion principle along lines Smith had theorized, and the project's own final reports, also since declassified, describe no confirmed detection of the sought anomalies and no experimental support for Smith's underlying antigravity-adjacent theoretical framework, leading to the project's quiet discontinuation, a specific negative experimental result from Smith's own dedicated, multi-year research program that receives considerably less attention in popular retellings than the earlier, unverified 1950 memo does.

\section{Why the Memo Outlived the Program's Actual Findings}

The 1950 memo's dramatic characterization, a program more classified than the hydrogen bomb, is a considerably more compelling and easily excerpted claim than Project Magnet's own more mundane and ultimately negative multi-year experimental conclusion, a pattern, an exciting unverified claim propagating further than its own program's later, quieter correction, that recurs across several cases addressed elsewhere in this series, including the Hendershot generator's early press coverage relative to its later technical correction, and that here specifically involves the same individual's own later research failing to substantiate the premise his earlier, less rigorously sourced memo had advanced.

\section{The Entropic Gravity Framing}

Smith's theoretical framework, proposing that a rotating or oscillating magnetic field configuration could couple to and manipulate gravitational binding, shares the general absence-of-coupling-mechanism problem addressed throughout this series regarding magnetic field based antigravity claims, including the Searl Effect Generator and Hamel spinning wheel entries, and Project Magnet's own negative experimental findings are consistent with, rather than an exception to, this general pattern of magnetic-field-based gravity claims failing to produce confirmed detection when actually, dedicatedly tested.

\section{Objections}

Supporters argue that the Canadian government's willingness to fund a multi-year dedicated research observatory demonstrates that Smith's underlying claims were taken seriously at an official level, lending them credibility beyond a typical unverified fringe claim. Government funding of an exploratory research program, as addressed in the Boeing GRASP entry elsewhere in this series regarding a different claim, demonstrates that the claim was judged worth a bounded investigation, not that the investigation confirmed the claim, and Project Magnet's actual negative findings, produced by that same government-funded investigation, are the more directly relevant evidence regarding the claim's validity than the funding decision that preceded them.

\section{Conclusion}

Smith's widely circulated 1950 memo documents his own secondhand, uncorroborated account of a private conversation rather than any independently verified government finding, and his own subsequent multi-year Project Magnet investigation, specifically designed to detect the phenomena and physics the memo described, concluded without confirming either, a negative result from the claim's own dedicated research program that its continued popular citation routinely omits.

\end{document}
